\documentclass[10pt,a4paper]{article}
\usepackage[margin=19mm]{geometry}
\usepackage{amsmath,amssymb,microtype,hyperref}
\setlength{\parindent}{0pt}
\setlength{\parskip}{5pt}
\begin{document}
\begin{center}
{\Large Neighbourhood enrichment: distance, background and prior}\par
{\small Historical correspondence and proposed simpler score; single-chain contract}
\end{center}
\textbf{Question and counting units.}
For a query receptor $q$ and epitope $X$, ask whether compatible VDJdb receptors
are more frequent near $q$ than receptors from a declared background.
A junction includes its conserved Cys and terminal Phe/Trp anchors.
Let $R_X$ be the eligible reference population and $B$ the eligible control population;
$N=|R_X|$ and $M=|B|$ are their sizes.
The ratios below require $N>0$ and $M>0$; empty populations are unavailable,
not zero-enrichment observations.
Use the same observation unit, distance and exact-identity exclusion for both.
Eligibility and exclusions must precede denominator calculation in the proposed model.
Historical denominators are preserved separately for reproduction.

\textbf{One additive distance.}
For amino acids $a,b$, let $s(a,b)$ be their BLOSUM62 entry and define
\[
c(a,b)=\begin{cases}0,&a=b,\\ \max(0,\min(4,4-s(a,b))),&a\ne b.\end{cases}
\]
This is the clipped substitution cost used by TCRdist3, not the BLOSUM Gram penalty.
Let $\ell_q,\ell_r$ denote junction lengths, $h=|\ell_q-\ell_r|$, and
$L=\max(\ell_q,\ell_r)$.
An allowed alignment $A$ has one contiguous length-difference block;
$(i,j)\in A$ denotes an aligned residue pair, indexed from zero.
With nonnegative positional weights $w_{L,k}$ and gap charge $g(h)$, use
\begin{equation}
\boxed{d(q,r)=d_V(q,r)+
\min_A\left[3\sum_{(i,j)\in A}w_{L,\max(i,j)}c(q_i,r_j)+g(h)\right].}
\end{equation}
Here $d_V$ is the sum of aligned CDR1, CDR2 and CDR2.5 substitution costs
from the two V-allele templates; each loop has weight one.
Missing or unresolved V calls do not imply zero loop distance.

\textbf{Flanks versus the junction core.}
TCRdist3 uses unit interior weights and excludes the first three and last two
junction residues. The author-requested comparison excludes three at \emph{both} ends.
Keep full strings for exact-identity exclusion even when end weights are zero:
different junctions with the same trimmed core are still distinct receptors.
An alternative is the shipped end-anchored empirical profile:
\[
u_{L,k}=1-\max\{p_V(k),p_J(L-1-k)\},\qquad
w_{L,k}=u_{L,k}\big/\left(L^{-1}\sum_{t=0}^{L-1}u_{L,t}\right).
\]
Here $p_V$ and $p_J$ are measured probabilities that a one-substitution
neighbour shares the epitope, indexed from the V and J ends; distant offsets
use the last pooled profile value. This discounts germline-rich flanks and
emphasises the central region. It is not an exact V/NDN/J boundary assignment.
A separate existing gene-specific retention profile estimates germline origin;
using it requires a symmetric query/reference weighting rule.
The ends can still carry gene information, which is represented explicitly by $d_V$.

\textbf{What does $g$ mean?}
It is a distance penalty, not a similarity reward.
For TCRdist3 repertoire defaults, $g(h)=12h$: the inner gap cost is four per
residue and the entire junction distance has outer weight three.
The gap position is optimised over the restricted central interval.
This is neither Smith--Waterman nor general multiple-indel alignment.
An affine alternative has $g(0)=0$ and
$g(h)=g_{\rm open}+(h-1)g_{\rm extend}$ for $h>0$;
linear is the special case $g_{\rm open}=g_{\rm extend}$.
A smoothly varying alternative charge is a different model, not a TCRdist3 default.
No discontinuous 20-residue route switch is required.

\newpage
\begin{center}{\large A count ratio with an explicit background}\end{center}
\textbf{Observed and expected neighbours.}
Fix a radius $D$ in the units of $d$ before counting.
It is a radius, not a diameter.
Write $n$ for the nonexact reference count within $D$, and $m$ for the
nonexact control count under the identical predicate.
If background ball probability is $p_B$, a sample of $N$ independent receptors
from that background has expected count $Np_B$.
Thus the basic enrichment and its finite-control ranking version are
\begin{equation}
E=\frac{n/N}{m/M},\qquad
\boxed{E_J=\frac{n}{N\widetilde p_B},\quad
\widetilde p_B=\frac{m+1/2}{M+1}.}
\end{equation}
The second expression uses the posterior-predictive probability under a
Jeffreys $\mathrm{Beta}(1/2,1/2)$ prior for the binary background-ball event.
It leaves $n=0$ at zero and avoids an arbitrary expected-count floor.
It is a regularised enrichment ranking, not a Bayes factor or a calibrated P-value.
Always report $n,N,m,M,N\widetilde p_B$ and zero-control support.
For $m/M>0$ and increasing $M$, $E_J$ approaches $E$;
with no control hits, its value depends on the available control size.
If both populations share the same distribution, their unsmoothed population
frequency ratio is one. Searching radii or selecting the best route requires
separate inferential calibration; a fixed-ball test cannot be reused unchanged.

\textbf{Three backgrounds, three questions.}
A generative V(D)J background measures enrichment beyond recombination accessibility.
A real control repertoire measures enrichment beyond its observed biological distribution;
healthy does not establish naive-cell provenance.
VDJdb excluding $X$ measures separation from other annotated specificities and
inherits database ascertainment bias.
They are separate named comparisons, not interchangeable estimates of the same null.
Sequence $P_{\rm gen}(q)$ is not ball probability $p_B$;
$P_{\rm gen}$ stratification diagnoses whether one method favours common rearrangements.
V-loop-inclusive counts require retained V annotations in controls.
Current historical control loading keeps only unique junctions and discards V/J.

\textbf{A prior without double counting.}
If a categorical V/J prior is wanted instead of embedding that information in
$d$, let $z$ denote one declared category and $N_z,M_z$ its reference/control sizes.
Let $n_z,m_z$ be geometric counts within those conditional populations.
Then the exact unsmoothed factorisation is
\[
\frac{N_z/N}{M_z/M}\;\frac{n_z/N_z}{m_z/M_z}
=\frac{n_z/N}{m_z/M}.
\]
The first factor is categorical enrichment; the second is conditional geometric
enrichment. Multiplying an unconditioned V-loop score by another V prior counts
related evidence twice. Empty or unavailable strata require explicit dispositions.

\textbf{Historical model and status.}
Original vdjmatch sums
$a(q,r)e^{-p(q,r)/400}/\max\{N\widehat F_q^{\rm edit}(H(q,r)),0.01\}$
over nonexact equal-length neighbours with at most five substitutions.
Here $a$ is the historical V factor, $p$ the positional BLOSUM Gram penalty,
$H$ the substitution count, and $\widehat F^{\rm edit}$ the control edit-ball mass.
The exponential is a historical attenuation heuristic, not a statistical derivation.
The matched-PSSM experiment substitutes the same positional geometry in that
background mass, but retains attenuation and, when the original prior is
requested and $N<500$, the historical
$\log(S+10^{-6})+G$ fusion, where $S$ is the summed density and $G$ the original
sum of marginal gene/length log-frequency ratios.
Otherwise the output remains $S$, including when $N\ge500$.
Its affine gap is $2800+1400(h-1)$ in different integer units; it is not $12h$.
Neither kernel choice passed all-task nonregression.
Equation (2) is the proposed simpler normalisation, not an accepted replacement.
Old scores remain reproducible; paired joint calibration and M3 are outside this note.
\end{document}
